\documentclass[12pt]{article}
\usepackage[utf8]{inputenc}
\usepackage[T1]{fontenc}
\usepackage[brazil]{babel}
\usepackage{geometry}
\usepackage{xcolor}
\usepackage{enumitem}
\usepackage{tikz}
\usetikzlibrary{positioning, arrows.meta, shapes.geometric}

% Definição da cor azul do título (aproximada da imagem)
\definecolor{meuazul}{RGB}{31, 108, 173}

% Configurações de margem
\geometry{a4paper, margin=2.5cm}

\begin{document}

\begin{center}
    \textcolor{meuazul}{\huge \textbf{Lista de Exercícios 02: Segurança da Informação (Parte 2)}}
\end{center}

\vspace{1cm}

% Configuração da lista principal (Questões)
\begin{enumerate}[
    label=\textbf{Questão \arabic*:},
    leftmargin=2.8cm, % Recuo para o texto alinhar sob o título da questão
    labelsep=0.3cm,   % Espaço entre "Questão X:" e o título
    align=left,
    labelwidth=2.5cm  % Largura fixa para o rótulo
]

    % Questão 1
    \item \textbf{Definição e Objetivos da Criptografia} \\
    A criptografia é uma ferramenta fundamental na engenharia de segurança. Explique, com suas palavras, o que é o processo de criptografia e quais são seus principais objetivos no contexto da proteção de dados.

    \vspace{0.4cm}

    % Questão 2
    \item \textbf{Componentes do Processo Criptográfico} \\
    Diferencie os seguintes termos fundamentais:
    \begin{enumerate}[label=\alph*), leftmargin=0.8cm, topsep=2pt]
        \item Texto Simples (\textit{Plaintext});
        \item Algoritmo Criptográfico;
        \item Texto Cifrado (\textit{Ciphertext});
        \item Chave (\textit{Key}).
    \end{enumerate}

    \vspace{0.4cm}

    % Questão 3
    \item \textbf{Criptografia Simétrica} \\
    A criptografia simétrica, ou de chave privada, é amplamente utilizada devido à sua eficiência.
    \begin{enumerate}[label=\alph*), leftmargin=0.8cm, topsep=2pt]
        \item Qual é a principal característica do uso de chaves nesse método?
        \item Cite uma vantagem e a principal desvantagem (desafio) da criptografia simétrica.
    \end{enumerate}

    \vspace{0.4cm}

    % Questão 4
    \item \textbf{Criptografia Assimétrica} \\
    Diferente da simétrica, a criptografia assimétrica utiliza um par de chaves (pública e privada).
    \begin{enumerate}[label=\alph*), leftmargin=0.8cm, topsep=2pt]
        \item Explique a função de cada uma das chaves no processo de envio de uma mensagem segura.
        \item Por que a criptografia assimétrica costuma ser combinada com a simétrica em aplicações reais (como no protocolo HTTPS)?
    \end{enumerate}

\end{enumerate}

\vspace{1cm}
\hrule
\vspace{0.5cm}

\begin{center}
    \textcolor{meuazul}{\Large \textbf{Respostas Sergio Novak}}
\end{center}

\vspace{0.5cm}

\begin{description}[style=nextline, font=\bfseries\color{meuazul}]
    
    \item[Resposta 1:] A criptografia, no contexto da engenharia de segurança, é um mecanismo de defesa essencial para proteger os \textbf{ativos do sistema}. Seu objetivo não é apenas o segredo, mas a garantia de atributos de dependabilidade como a \textbf{Confidencialidade} (impedir acesso não autorizado) e a \textbf{Integridade} (garantir que os dados não foram corrompidos).

    \begin{center}
        \begin{tikzpicture}[node distance=1cm, auto, >=Stealth, thick]
            \node (P1) [draw, rounded corners, fill=white] {Texto Simples};
            \node (E) [draw, fill=meuazul!15, right=of P1] {Cifragem};
            \node (C) [draw, fill=gray!10, right=of E] {Texto Cifrado};
            \node (D) [draw, fill=meuazul!15, right=of C] {Decifragem};
            \node (P2) [draw, rounded corners, fill=white, right=of D] {Texto Simples};
            
            \draw [->] (P1) -- (E);
            \draw [->] (E) -- (C);
            \draw [->] (C) -- (D);
            \draw [->] (D) -- (P2);
            
            \node (K) [above=0.5cm of E] {Chave};
            \draw [->, dashed] (K) -- (E);
            \node (K2) [above=0.5cm of D] {Chave};
            \draw [->, dashed] (K2) -- (D);
        \end{tikzpicture}
    \end{center}

    \item[Resposta 2:] Para sistemas resilientes, distinguimos: 
    \begin{itemize}[label=--]
        \item \textit{Plaintext}: O dado vulnerável original.
        \item \textit{Algoritmo}: A lógica de transformação (deve ser pública e robusta).
        \item \textit{Ciphertext}: O dado protegido para trânsito em meios inseguros.
        \item \textit{Chave}: O parâmetro crítico de segurança que deve ser gerido com rigor.
    \end{itemize}

    \item[Resposta 3:] A característica principal da criptografia simétrica é o compartilhamento de uma \textbf{única chave secreta}. 
    \begin{itemize}[label=--]
        \item \textit{Vantagem}: Alta eficiência computacional, ideal para grandes volumes de dados.
        \item \textit{Desafio}: O problema socio-técnico do gerenciamento e distribuição segura das chaves entre as partes.
    \end{itemize}

    \begin{center}
        \begin{tikzpicture}[>=Stealth, thick, node distance=2cm]
            \node (UserA) [draw, circle, fill=meuazul!10] {A};
            \node (UserB) [draw, circle, fill=meuazul!10, right=4cm of UserA] {B};
            \draw [<->] (UserA) -- node[above] {Mesma Chave (K)} (UserB);
            \node at (3,0.5) {\textbf{Simétrica}};
        \end{tikzpicture}
    \end{center}

    \item[Resposta 4:] A criptografia assimétrica utiliza pares de chaves: a pública para cifragem/verificação e a privada para decifragem/assinatura. 
    Em aplicações como o \textbf{HTTPS}, utilizamos um modelo híbrido: a criptografia assimétrica resolve o problema da troca segura de chaves (gerenciamento), e a criptografia simétrica assume a sessão para manter a performance do sistema.

    \begin{center}
        \begin{tikzpicture}[>=Stealth, thick]
            \node (Client) [draw, rectangle, fill=white] {Cliente};
            \node (Server) [draw, rectangle, fill=white, right=5cm of Client] {Servidor};
            
            \draw [->] (Client) -- node[above, font=\tiny] {1. Solicita conexão} (Server);
            \draw [<-] (Client) -- node[below, font=\tiny, pos=0.5, yshift=-0.2cm] {2. Envia Chave Pública} (Server);
            \draw [->] (Client) -- node[above, font=\tiny, sloped] {3. Cifra Chave de Sessão} (Server);
            \draw [<->, color=meuazul, line width=1pt] (Client) to[bend left=20] node[above, font=\bfseries\tiny] {Canal Seguro (Simétrico)} (Server);
        \end{tikzpicture}
    \end{center}
\end{description}

\end{document}
